\section{Направления дальнейшего развития}

Разработанный прототип можно развивать
в нескольких направлениях. Одно из них связано с расширением признакового
пространства. В текущей версии система опирается на историю культур, базовые характеристики поля и региональный контекст.
Добавление почвенных, климатических, экономических и производственных признаков позволило бы приблизить рекомендации к
более реалистичному сценарию поддержки принятия решений.

Отдельным направлением является адаптация предложенного подхода к российскому сельскохозяйственному контексту. Для
этого необходимы локальные данные о культурах на уровне полей и годов, региональные агрономические правила, а также
подключение почвенных, климатических и производственных признаков. Такая адаптация позволила бы проверить
переносимость разработанной методологии и уточнить интерпретационный слой с учётом особенностей конкретных регионов.

Ещё одно направление — более детальный анализ уверенности. В работе были рассмотрены группы уверенности и
диагностическая оценка калибровки вероятностей. В дальнейшем этот анализ можно расширить: отдельно изучать калибровку по
классам культур, регионам, типам последовательностей и сложным случаям, где первая рекомендация модели оказывается
неверной, но правильная культура сохраняется среди ближайших альтернатив.

Отдельно можно развивать граф знаний. Графовый слой может быть дополнен новыми типами связей,
литературными правилами и экспертными корректировками. Это позволит сделать интерпретационный контекст более содержательным
и полезным для анализа отдельных кандидатов.

Отдельным направлением является развитие пространственного интерфейса. В текущей работе карта используется как
демонстрация сценария анализа результатов на уровне полей. В дальнейшем такой интерфейс можно расширить: объединить карту,
список рекомендаций, оценку уверенности и пояснения графа в едином пользовательском сценарии.

Таким образом, дальнейшее развитие работы связано не только с улучшением предиктивной модели, но и с расширением всей
системы поддержки принятия решений: данных, диагностического анализа уверенности, интерпретационного слоя, интерфейса
анализа и адаптации подхода к российскому сельскохозяйственному контексту.
